% U5 WS2 -- integrated rate laws and half-life. Block 3.
\documentclass{apchem-worksheet}

\apcunit{5}
\apcunittitle{Kinetics}
\apcdoctitle{Worksheet 2 \textbullet\ Concentration Over Time}
\apcsubtitle{CED 5.3 \textbullet\ say \emph{which} plot is linear, not just ``the graph''}

\begin{document}
\wsheader[Zero: {$[\ce{A}] = -kt + [\ce{A}]_0$}, {$t_{1/2} = [\ce{A}]_0/2k$}
\textbullet\ First: {$\ln[\ce{A}] = -kt + \ln[\ce{A}]_0$},
{$t_{1/2} = 0.693/k$} \textbullet\ Second:
{$1/[\ce{A}] = kt + 1/[\ce{A}]_0$}, {$t_{1/2} = 1/(k[\ce{A}]_0)$}.]

\problem[]{\ced{5.3} Complete the table from memory, then check it against
the data strip:}
\begin{parts}
  \item Zero order --- linear plot: \answerblank[3.4cm]{$[\ce{A}]$ vs $t$},
        slope \answerblank[1.6cm]{$-k$}
  \item First order --- linear plot:
        \answerblank[3.4cm]{$\ln[\ce{A}]$ vs $t$}, slope
        \answerblank[1.6cm]{$-k$}
  \item Second order --- linear plot:
        \answerblank[3.4cm]{$1/[\ce{A}]$ vs $t$}, slope
        \answerblank[1.6cm]{$+k$}
\end{parts}

\problem[]{\ced{5.3} A first-order reaction has
$k = \SI{0.0250}{\per\second}$ and $[\ce{A}]_0 = \SI{0.500}{\Molar}$.}
\begin{parts}
  \item Find $[\ce{A}]$ after \SI{60.0}{\second}. \workspace[2.4cm]
        \keyonly{\begin{apcanswer}$\ln[\ce{A}] = -(0.0250)(60.0) +
        \ln(0.500) = -1.50 - 0.693 = -2.193$;
        $[\ce{A}] = e^{-2.193} =
        \mathbf{0.112}~\si{\Molar}$\end{apcanswer}}
  \item Find the half-life. \answerblank[2.6cm]{\SI{27.7}{\second}}
  \item After two half-lives, what fraction remains, and does your answer
        to (a) look consistent? \answerlines[2]{One quarter, so
        \SI{0.125}{\Molar} at \SI{55.4}{\second}. The value at
        \SI{60}{\second} is slightly lower, \SI{0.112}{\Molar}, which is
        exactly what it should be.}
\end{parts}

\problem[]{\ced{5.3} A second-order reaction has
$k = \SI{0.15}{\per\Molar\per\second}$ and
$[\ce{A}]_0 = \SI{0.20}{\Molar}$.}
\begin{parts}
  \item Find $[\ce{A}]$ after \SI{100}{\second}. \workspace[2.2cm]
        \keyonly{\begin{apcanswer}$1/[\ce{A}] = (0.15)(100) + 1/0.20 =
        15 + 5.0 = 20$; $[\ce{A}] =
        \mathbf{0.050}~\si{\Molar}$\end{apcanswer}}
  \item Find the half-life. \answerblank[2.6cm]{\SI{33}{\second}}
  \item After the first half-life, will the second take the same time,
        less, or more? Justify. \answerlines[3]{More. For second order
        $t_{1/2} = 1/(k[\ce{A}]_0)$ depends on the starting concentration,
        and once half has reacted the remaining concentration is lower, so
        the next half-life is twice as long.}
\end{parts}

\problem[]{\ced{5.3} A zero-order reaction has
$k = \SI{0.020}{\Molar\per\second}$ and $[\ce{A}]_0 = \SI{1.00}{\Molar}$.}
\begin{parts}
  \item Find $t_{1/2}$. \answerblank[2.4cm]{\SI{25}{\second}}
  \item Find $[\ce{A}]$ after \SI{30}{\second}.
        \answerblank[2.4cm]{\SI{0.40}{\Molar}}
  \item At what time does the reactant run out completely? What does that
        tell you about zero-order behaviour?
        \answerlines[3]{At \SI{50}{\second}, since the concentration falls
        linearly. Zero-order kinetics cannot continue past that point --- a
        real reaction stops being zero order once the reactant becomes
        scarce, because the mechanism causing the saturation (a full
        catalyst surface, say) can no longer be maintained.}
\end{parts}

\problem[]{\ced{5.3} Identify the order from each observation:}
\begin{parts}
  \item Half-life is \SI{25.0}{\minute} at the start and
        \SI{25.0}{\minute} again later: \answerblank[3cm]{first order}
  \item $k$ for that reaction: \answerblank[3cm]{\SI{0.0277}{\per\minute}}
  \item A plot of $1/[\ce{A}]$ vs $t$ is a straight line with slope
        $+0.42$: \answerblank[6.4cm]{second order, $k = 0.42$
        \si{\per\Molar\per\second}}
  \item A plot of $[\ce{A}]$ vs $t$ is a straight line:
        \answerblank[3cm]{zero order}
\end{parts}

\problem[]{\ced{5.3} A student writes ``the graph was linear, so the
reaction is first order.'' Explain why this earns no credit and write the
sentence that would.
\answerlines[3]{All three orders produce \emph{some} linear plot --- the
statement does not distinguish them. The answer must name the quantity
plotted: ``the plot of $\ln[\ce{A}]$ against time is linear, therefore the
reaction is first order in A.''}}

\problem[]{\ced{5.3} Why is a constant half-life such strong evidence of
first-order kinetics?
\answerlines[3]{Of the three expressions, only $t_{1/2} = 0.693/k$ contains
no $[\ce{A}]_0$. For zero order the half-life shortens as the reaction
proceeds and for second order it lengthens, so a half-life that does not
change with concentration is a signature of first order.}}

\problem[]{\ced{5.3} A radioactive isotope has a half-life of
\SI{8.0}{\day}. A sample starts at \SI{40.0}{\milli\gram}.}
\begin{parts}
  \item Mass remaining after \SI{24}{\day}:
        \answerblank[2.6cm]{\SI{5.0}{\milli\gram}}
  \item How did you get it without using the integrated rate law?
        \answerlines[2]{\SI{24}{\day} is exactly three half-lives, so the
        mass halves three times: $40.0 \to 20.0 \to 10.0 \to 5.0$~mg.}
  \item Radioactive decay is always first order. What does that imply about
        the time to go from 40.0 to 20.0 mg versus 20.0 to 10.0 mg?
        \answerblank[4.4cm]{the same --- \SI{8.0}{\day} each}
\end{parts}

\begin{spiralstrip}{Unit 5 \textbullet\ blocks 1--2}
  \item Orders come from: \answerblank[2.6cm]{experiment}
  \item Units of $k$, first order: \answerblank[2.4cm]{\si{\per\second}}
  \item Doubling [A] quadruples the rate; order in A?
        \answerblank[1.6cm]{2}
  \item Rate for $\ce{2A -> B}$ in terms of A:
        \answerblank[4cm]{$-\tfrac{1}{2}\Delta[\ce{A}]/\Delta t$}
  \item Initial rate is measured at: \answerblank[2cm]{$t=0$}
\end{spiralstrip}

\end{document}
